\section{Estimation and inference}\label{sec:inference}

\paragraph{Sampling.} We observe an i.i.d. pooled sample $O_i=(Z_i,X_i,(1-Z_i)Y_i)$, $i=1,\dots,n$, with $\PP(Z=1)=\pi$, $X\mid Z=1\sim P$ and $X\mid Z=0\sim Q$; for $\theta$ only the outcomes of donor units are needed. When a unit (a customer) generates several observations (calls), $O_i$ is the whole block of that unit and everything below applies at the unit level. We assume that the channel and target-group membership are properties of the unit, constant across its observations, and that $P$, $Q$, $m_0$ and $\theta$ are defined at the observation level, so that each unit weighs according to its number of observations. If that number depends on $Y(0)$ given $X$, the unit-level and observation-level estimands, and the corresponding versions of (A2), differ.

\paragraph{Estimators.} With $\hat o=\hat p/(1-\hat p)$ and $n_1$ the number of target-group observations,
\[
\hat\theta_{\mathrm{H}}=\frac{\sum_{i:Z_i=0}\hat o(X_i)\,Y_i}{\sum_{i:Z_i=0}\hat o(X_i)},\qquad
\hat\theta_{\mathrm{DR}}=\frac{\sum_{i:Z_i=0}\hat o(X_i)\,(Y_i-\hat m(X_i))}{\sum_{i:Z_i=0}\hat o(X_i)}+\frac1{n_1}\sum_{i:Z_i=1}\hat m(X_i).
\]
The first is the Hájek weighting estimator and the second the sample version of $\psi$ in Proposition~\ref{prop:dr} with the weighting term normalized. Both use cross-fitting: units are split into $K$ folds and, for observations in fold $k$, $\hat p$ and $\hat m$ are trained on the other folds. With a known cutoff, $\hat o$ is replaced by $\ind\{s\le c\}\,\hat e/(1-\hat e)$ (Proposition~\ref{prop:cutoff}), and the sample also includes the premium units outside the target group. Under population sampling of $(D,S,X,Y)$, the corresponding influence function is $\PP(S=1)^{-1}\bigl[S(m_0-\theta)+(1-D)\,\ind\{s\le c\}\,e/(1-e)\,(Y-m_0)\bigr]$; we expect an analogue of Theorem~\ref{thm:normality} but do not prove it.

\subsection{Weight trimming}\label{sec:trimming}

Trimming the weights changes the estimand. For $\tau>0$, let $w_\tau=\min(w,\tau)$ and $\varphi_\tau(t)=\min(1,\tau/t)$ for $t>0$. With the true weights, the trimmed self-normalized estimator converges to
\[
\theta_\tau=\frac{\E[w_\tau(X)\,Y\mid D=0]}{\E[w_\tau(X)\mid D=0]}.
\]

\begin{proposition}[Trimming bias]\label{prop:trimming}
Under (A1)--(A4):
\begin{enumerate}
  \item $\displaystyle\theta_\tau-\theta=\frac{\Cov_P\bigl(m_0(X),\varphi_\tau(w(X))\bigr)}{\E_P[\varphi_\tau(w(X))]}$.
  \item If $g(t)=\E_P[m_0(X)\mid w(X)=t]$ has a nonincreasing version, $\theta_\tau\ge\theta$; if it has a nondecreasing version, $\theta_\tau\le\theta$.
  \item If $Y\ge0$, the unnormalized version is always below: $\E[w_\tau(X)\,Y\mid D=0]\le\theta$.
\end{enumerate}
\end{proposition}
\begin{proof}
Since $P(\{w=0\})=\int_{\{w=0\}}w\dd Q=0$, on $\{w>0\}$ we have $w_\tau=\varphi_\tau(w)\,w$. By the tower property and the change of measure, $\E[w_\tau(X)Y\mid D=0]=\int w_\tau m_0\dd Q=\E_P[\varphi_\tau(w)\,m_0]$ and $\E[w_\tau(X)\mid D=0]=\E_P[\varphi_\tau(w)]\in(0,1]$; subtracting $\theta=\E_P[m_0]$ gives (1). For (2), the tower property gives $\Cov_P(m_0,\varphi_\tau(w))=\Cov_P(g(T),\varphi_\tau(T))$ with $T=w(X)$. The function $\varphi_\tau$ is nonincreasing; if $g$ is too, $(g(T)-g(T'))(\varphi_\tau(T)-\varphi_\tau(T'))\ge0$ for an independent copy $T'$, and its expectation is $2\Cov_P(g(T),\varphi_\tau(T))$ (Chebyshev's association inequality); the nondecreasing case is symmetric. For (3), $\varphi_\tau\le1$ and $m_0\ge0$ give $\E_P[\varphi_\tau(w)m_0]\le\theta$.
\end{proof}

\paragraph{Interpretation.} A large $w(x)$ means that profile $x$ is much more frequent in the target group than in the donor channel. There is no default direction: it depends on whether, within the target group, the profiles most typical of the premium channel respond more or less in the basic channel. If they respond less, $g$ is decreasing and trimming overestimates $\theta$ and hence underestimates the loss; in our simulation they respond more (Section~\ref{sec:res-trimming}). The condition concerns $g$ as a function of the weight, not of $X$, and it can be checked approximately, because $m_0$ is identified on the support.

\paragraph{Example~\ref{ex:toy} (trimming).} With $\tau=2$, the weights $(0.2,\,1,\,3)$ become $(0.2,\,1,\,2)$ and $\theta_\tau=1.2\,\%/0.8=1.5\,\%$, against $\theta=1.3\,\%$. The formula agrees: $\varphi_\tau=(1,\,1,\,2/3)$, $\E_P[\varphi_\tau]=0.8$ and the covariance is $0.16$ points, so the bias is $0.2$ points. The unnormalized version gives $1.2\,\%$, below $\theta$.

\begin{corollary}[Trimming with the doubly robust estimator]\label{cor:drtrimming}
By Proposition~\ref{prop:dr} with $\bar w=w_\tau$, $\psi(w_\tau,\bar m)-\theta=-\int_{\{w>\tau\}}(w-\tau)(m_0-\bar m)\dd Q$.
\end{corollary}
The corollary concerns the unnormalized form $\psi$: trimming biases only to the extent that the outcome model is wrong in the trimmed region. For the normalized form we use, the bias is $\E_Q[w_\tau(m_0-\bar m)]/\E_Q[w_\tau]-\E_P[m_0-\bar m]$, which vanishes if $\bar m=m_0$ on the whole support, but not necessarily if the model is right only in the trimmed region. Since Proposition~\ref{prop:dr} holds for any weights, when estimated weights are trimmed the bias depends on both errors, and it vanishes if $\bar m=m_0$.

\paragraph{Estimated weights.} In practice $\hat w$ is trimmed, and trimming removes both true weight and estimation noise, so the sign in Proposition~\ref{prop:trimming} is not guaranteed. In the companion note, included in the repository that accompanies this paper, we prove that it is preserved if the estimation error is a profile-independent multiplicative noise with a log-concave density. Before trimming, the effect of trimming computed with $\hat m$ instead of $Y$, or $\hat m$ by mass deciles of $\hat w$, indicates which way it will push. \citet{crump2009} trim on the propensity and redefine the estimand; \citet{ma2020} keep the estimand and characterize the bias of trimming by discarding in the unnormalized estimator, of which item 3 is the version for capping; the contribution of Proposition~\ref{prop:trimming} is the normalized case, with its exact formula. \citet{liu2023} treat truncation as a change of estimand for the effect on the treated.

\subsection{Variance and asymptotic normality}

\paragraph{Influence function.} Let $o=p/(1-p)$ and $\sigma^2(x)=\Var(Y\mid X=x,D=0)$. The influence function of the doubly robust estimator, of the same form as the efficient one for the counterfactual mean of the treated \citep{hahn1998}, is
\begin{gather*}
\phi(O)=\frac1\pi\Bigl[Z\,(m_0(X)-\theta)+(1-Z)\,o(X)\,(Y-m_0(X))\Bigr],\\
V=\E[\phi^2]=\frac{\Var_P(m_0(X))}{\pi}+\frac{\E_Q[w(X)^2\sigma^2(X)]}{1-\pi},
\end{gather*}
where the cross term vanishes because $Z(1-Z)=0$ and the second term uses $o=\pi w/(1-\pi)$ (Theorem~\ref{thm:odds}). This closed form is for i.i.d. observations; with several observations per unit, see the appendix.

\begin{proposition}[When the variance is finite]\label{prop:variance}
If $Y$ is bounded and $0<\sigma_-^2\le\sigma^2\le\sigma_+^2$ $Q$-almost surely, then $V<\infty$ if and only if $\E_Q[w^2]<\infty$, that is, if and only if $\chi^2(P\|Q)=\E_Q[w^2]-1<\infty$. Absolute continuity $P\ll Q$ does not imply this condition.
\end{proposition}
\begin{proof}
The first term of $V$ is finite because $m_0$ is bounded; the second lies between $\sigma_-^2\,\E_Q[w^2]/(1-\pi)$ and $\sigma_+^2\,\E_Q[w^2]/(1-\pi)$. To see that $P\ll Q$ is not enough, let $Q$ be uniform on $(0,1)$ and $w(x)=x^{-1/2}/2$, which integrates to 1; then $P\ll Q$ but $\E_Q[w^2]=\int_0^1(4x)^{-1}dx=\infty$.
\end{proof}

Thus (A3) suffices to identify $\theta$ but not to estimate it at rate $n^{-1/2}$: with the conditional variance bounded below, one also needs $\chi^2(P\|Q)<\infty$, and standard theory requires something stronger, bounded weights.

\paragraph{Effective sample size.} Kish's $\neff=(\sum_i\hat w_i)^2/\sum_i\hat w_i^2$ over the $n_0$ donor observations satisfies, with the true weights, $\neff/n_0\to1/\E_Q[w^2]=1/(1+\chi^2(P\|Q))$, so it estimates the $\chi^2$ divergence; with estimated weights it can move either way (noise in $\hat w$ lowers it, shrinkage by a regularized classifier raises it), and in our simulation it falls below the theoretical value except under low overlap, where the theoretical value is itself imprecise (Section~\ref{sec:design}). In finite samples it is always finite, even when $\chi^2=\infty$, and it then goes to zero very slowly and noisily: in the counterexample of Proposition~\ref{prop:variance}, the median over replications goes from 0.45 with $n_0=10^3$ to about 0.22 with $n_0=10^7$, and with $n_0=10^3$ the 1st and 99th percentiles across replications are 0.05 and 0.65. A reasonable $\neff$ does not guarantee that $\chi^2$ is finite; one should also examine the stability of the upper quantiles of $\hat w$.

\begin{theorem}[Asymptotic normality of the doubly robust estimator]\label{thm:normality}
Assume (A1)--(A4), bounded $Y$ and, in addition:
\begin{enumerate}
  \item \emph{Strong overlap:} $p(X)\le1-\eta$ $Q$-almost surely for some $\eta>0$, that is, bounded $w$.
  \item Cross-fitting by unit with a fixed number $K$ of folds of size proportional to $n$, with $\hat p\le1-\eta$ and $\hat m$ bounded by the same bounds as $Y$.
  \item \emph{Rates:} $r_o=\max_k\|\hat o_k-o\|$ and $r_m=\max_k\|\hat m_k-m_0\|$ in $L^2(Q)$ go to 0 in probability, and $r_o\,r_m=o_P(n^{-1/2})$.
\end{enumerate}
Then $\sqrt n(\hat\theta_{\mathrm{DR}}-\theta)\to N(0,V)$ in distribution, and the sample mean of $\hat\phi(O_i)^2$ estimates $V$ consistently.
\end{theorem}
The proof, in Appendix~\ref{app:proof}, adapts Theorem 3.1 of \citet{chernozhukov2018} and its example of the effect on the treated: by Proposition~\ref{prop:dr}, the bias is bilinear in the errors of the auxiliary parameters (Neyman orthogonality), cross-fitting controls the empirical-process term without complexity conditions, and the denominator is handled by the delta method.

\paragraph{Beyond strong overlap.} Condition 1 is stronger than (A3) and than $\chi^2<\infty$. If the weights are unbounded but $\chi^2<\infty$, $V$ is finite, although the standard proof does not cover the case. If $\chi^2=\infty$ and the conditional variance is bounded below, no regular estimator converges at rate $\sqrt n$: this is the irregular identification case of \citet{khan2010}. \citet{ma2020} show that the limit may then be non-normal and propose subsampling inference, and \citet{ma2023} give adaptive inference for doubly robust estimators with trimming and bias correction. Our simulation lies in this regime: the channel index is unbounded inside the target region, so the weights are unbounded while $\chi^2$ is finite, and the implementation clips $\hat p$ to $[10^{-6},1-10^{-6}]$ rather than to $1-\eta$; its coverage results are evidence outside the conditions of Theorem~\ref{thm:normality}.

\paragraph{Intervals.} We use $\hat\theta_{\mathrm{DR}}\pm1.96\,(\hat V/n)^{1/2}$. Theorem~\ref{thm:normality} covers neither weight trimming nor the bootstrap, which we evaluate in the simulation (Sections~\ref{sec:res-trimming} and~\ref{sec:res-bootstrap}). If the target group is defined with an estimated threshold $\hat c$, its uncertainty adds a term to $V$ that we do not include: we treat the selection rule as fixed.
